%% ======================================================================
%% chap9/conclusion_main.tex
%%
%% Chapter 9: Conclusion and Future Work
%%
%% Closes the dissertation: summarises the four contribution categories
%% from Chapter 1, distils six principal findings, presents 11 threads of
%% future work (Table 9.1, with full descriptions in Appendix E), and
%% offers a closing reflection on the premise and its refinement.
%%
%% Sections:
%%   9.1 Summary of Contributions (enumerate list of 4)
%%   9.2 Principal Findings (enumerate list of 6)
%%   9.3 Future Work (Table 9.1 -- 11 threads)
%%   9.4 Closing Reflection
%%
%% External labels referenced:
%%   sec:contributions (1.4), chap:results (Ch 7),
%%   subsec:economic-crossover (8.2.2),
%%   app:future-work (Appendix E)
%%
%% Phantom labels embedded in Table 9.1 rows (referenced from earlier
%% chapters): subsec:future-work-external-validation (row 9.3.1),
%% subsec:bounded-vec-future (row 9.3.4),
%% subsec:audience-validation-future (row 9.3.5)
%%
%% Citations: none.
%% ======================================================================

\chapter{Conclusion and Future Work}\label{chap:conclusion}

This closing chapter restates the contributions, summarizes the principal findings, identifies the most important threads of future work, and offers a brief closing reflection.

\section{Summary of Contributions}\label{sec:summary-of-contributions}

Section~\ref{sec:contributions} introduced our four contribution categories. We summarize them briefly here:

\begin{enumerate}
\item \textbf{The artifact:} (pallets, runtime, ink! contract, node, benchmark, and simulation scripts) released under the MIT license as open source, intended as a starting point for further research.
\item \textbf{The formal documentation:} seven implementation specifications that address the pallet, XCM message flows, security threat model, royalty algorithm, precompile interface, NFT nesting model, and Snowbridge integration. These documents are intended as design knowledge rather than code documentation: they elucidate what the artifact accomplishes and the rationale behind each design decision.
\item \textbf{The empirical findings:} all 11 KPIs met their targets (Chapter~\ref{chap:results}); the Monte Carlo simulation produced the audience-dependent creator revenue result discussed in Section~\ref{subsec:economic-crossover}; the security audit yielded eight findings (one Medium and one Low remediated); and the Snowbridge end-to-end demonstration validates forward compatibility with Ethereum bridging.
\item \textbf{A worked example of DSR applied to a multi-component blockchain prototype:} the integration of pallet implementation, smart contract development, cross-chain interoperability, and economic simulation within a DSR framework is inadequately represented in the current body of DSR literature. Future researchers developing similar systems may consider this dissertation as a valuable reference template.
\end{enumerate}

\section{Principal Findings}\label{sec:principal-findings}

Six principal findings emerge from the work:

\begin{enumerate}
\item \textbf{The architecture works.} All planned functionality is implemented; all KPIs are met; all 56 unit tests pass on every commit.
\item \textbf{The major architectural decisions aged well.} The pallet-first decision was vindicated by the January 2026 ink! discontinuation; the RMRK 2.0 to \texttt{pallet-nfts} pivot, the unified rights token model, and the cross-chain operation pattern are likewise vindicated to varying degrees. In retrospect, we would not change any of the major decisions.
\item \textbf{The economic advantage is conditional on the creator's audience size.} CCRMS wins for creators with 10,000+ followers but loses to centralized platforms for creators with fewer than $\sim$2,000 followers because algorithmic discovery dominates. The headline number (CCRMS earns more on average) is true, but the conditional version is the more accurate framing for individual creator decisions.
\item \textbf{Cross-chain operation is operationally feasible but latency-bound.} XCM v3 is sufficient; cross-chain finality measured on a pre-Polkadot~2.0 toolchain serves as an upper bound, not a fundamental limit.
\item \textbf{Standards reimplementation is sometimes preferable to standards adoption} when the reference implementation is unmaintained; the RMRK 2.0 pivot is the exemplar.
\item \textbf{Forward compatibility is a defensive engineering goal.} Configuring the runtime to accept Ethereum-originated XCM messages at minimal additional cost in Phase~4 enabled the Snowbridge integration to succeed without pallet-level modifications.
\end{enumerate}

\section{Future Work}\label{sec:future-work}

11 threads of future work emerged from the research. Table~\ref{tab:future-work} summarizes each; Appendix~\ref{app:future-work} provides full descriptions.

\begin{longtable}{@{}L{0.06\linewidth} L{0.26\linewidth} L{0.14\linewidth} L{0.48\linewidth}@{}}
\caption{Future Work Summary}\label{tab:future-work}\\
\toprule
\textbf{\#} & \textbf{Thread} & \textbf{Priority} & \textbf{Key action} \\
\midrule
\endfirsthead

\multicolumn{4}{l}{\emph{Table~\ref{tab:future-work} continued from previous page}}\\
\toprule
\textbf{\#} & \textbf{Thread} & \textbf{Priority} & \textbf{Key action} \\
\midrule
\endhead

\midrule
\multicolumn{4}{r}{\emph{Continued on next page}}\\
\endfoot

\bottomrule
\endlastfoot

9.3.1\phantomsection\label{subsec:future-work-external-validation} & Production deployment & Near-term & Deploy to public testnet/mainnet; obtain third-party security audit \\
\addlinespace
9.3.2 & Polkadot~2.0 migration & Near-term & Target newer SDK release (no pallet changes); compress latency via Async Backing, Agile Coretime, Elastic Scaling \\
\addlinespace
9.3.3 & Smart contract layer migration & Near-term & Choose among Solidity rewrite, \texttt{wrevive} adoption, or contract removal (ink! discontinued Jan 2026) \\
\addlinespace
9.3.4\phantomsection\label{subsec:bounded-vec-future} & 50-child cap mitigation & Near-term & Raise bound + re-benchmark, replace with membership map, or introduce hierarchical nesting \\
\addlinespace
9.3.5\phantomsection\label{subsec:audience-validation-future} & User study & Near-term & Build minimal wallet/CLI; test with Web3 creators and consumers (highest-priority research extension) \\
\addlinespace
9.3.6 & Bridged-payment final hop & Near-term & Configure \texttt{PaymentCurrency} for foreign Ether or add swap step \\
\addlinespace
9.3.7 & Resale royalties & Medium-term & Add optional \texttt{resale\_royalty\_basis\_points}; requires on-chain marketplace for price verification \\
\addlinespace
9.3.8 & Capability-based delegation & Medium-term & Relax Finding~B remediation with time-limited delegate authorization \\
\addlinespace
9.3.9 & Monte Carlo extensions & Medium-term & Add creator psychology, network effects, longitudinal and regulatory modeling \\
\addlinespace
9.3.10 & Formal verification & Long-term & Apply Isabelle/HOL techniques to authorization model, royalty algorithm, cross-chain pattern \\
\addlinespace
9.3.11 & JAM transition & Long-term & Migrate to JAM when mainnet launches (mid-to-late 2026); no rebuild required \\
\end{longtable}

\section{Closing Reflection}\label{sec:closing-reflection}

CCRMS began with a premise: that creators deserve more of the value they generate and that blockchain technology offers a credible path to delivering it. We tested the premise by building a working prototype, measuring its operational properties, and simulating its economic outcomes. The premise held, but it was also \textbf{refined}: the Monte Carlo simulation contradicts the naive form (``decentralization is universally better for creators'') and supports the conditional form (``decentralization is better for creators who have outgrown the discovery value of centralized platforms''). An honest conditional is more valuable than a universal claim.

The Polkadot ecosystem changed substantially during the implementation window (Polkadot~2.0 moved to production, Snowbridge V2 went live, ink! was discontinued, JAM reached testnet, Story Protocol launched), and any of these could have invalidated the work under different architectural decisions. None did, because we chose conservatively: favoring maintained primitives, minimizing reliance on non-essential components, and preferring deployable patterns over speculative future features. In retrospect, the dissertation is a study in \textbf{defensive engineering applied to a fast-moving research domain}, and a case for \textbf{conservatively chosen, carefully documented architectural decisions as the most durable deliverable}.

CCRMS is a prototype in one corner of one blockchain ecosystem, addressing one aspect of one industry's structural problems. It is not the future of content rights management, and we do not claim it is. But it demonstrates that the future of content rights management can be built on open, decentralized, cryptographically verifiable infrastructure, and it provides a starting point for others to continue the work.
